% ---
% id: meat-core/frank-wolfe/rem:fo-general-format
% kind: rem
% title: "The first-order bound in the general format"
% summary: >-
%   The first-order theorem holds as written in the general format, since no
%   step of its proof reads a constraint row, the number of measured factors or
%   the number of Alice sectors. In trusted-source mode the linear minimisation
%   splits into one conic program per Alice sector.
% prominence: standard
% status: accepted
% uses: [thm:firstorder, def:first-order-setting, meat-core/thm:choi, meat-core/prop:a-neutral, meat-core/lem:branch-support]
% mentions: [cor:fo-dropped, lem:feasible-side-dual, def:relaxed-feasible-set, prop:relaxed-set-properties, rem:feas-tol, req:per-sector-coupling, rem:fo-interface, rem:fo-iterate-checks, rem:fo-returned-point, rem:fo-one-ruler, rem:fo-reachability, rem:fo-check-exposures, rem:fo-oracle-repair, rem:fo-iterate-repair, rem:fo-negative-values, meat-core/def:declaration]
% assumes: []
% refs: []
% provenance: newly-derived
% lean:
% ---

\begin{remark}\label{rem:fo-general-format}
\emph{(1) Theorem~\ref{thm:firstorder} is stated for a declaration in the
general format} of Definition~\ref{def:declaration} as well, with
$\Feas^{J}$ the set of lists over the joint sectors of
Theorem~\ref{thm:choi} and $J_s\coloneqq\bigoplus_{\mathsf t}
J_{(\mathsf t,s)}$ for every sector $s$ of the measured system, and the
last paragraph of its proof checks each step there. That check is made on
the general format itself, since the \tref{basic format} cannot express
the general case: a \tref{party-local} entry, a
\tref{transcript protocol} of more than two messages and an
\tref{ensemble form} measurement slot are each outside it, so no appeal to
``the same declaration with its factor tagged $A$ forgotten'' would carry
the general case, and none is made.
\emph{What Proposition~\ref{prop:a-neutral} does and does not contribute
here.} It says that restricting the variable to states block diagonal along
Alice's sectors costs nothing, which is why the joint-sector program is the
right program to run and not a weaker one; it is not what carries the
theorem, and the inclusion of the joint-sector Choi set into a Choi set of
the $A$-forgotten declaration is not used, nor is it available in general,
that declaration's Choi feasible set carrying only
\tref[lifted operator]{lifted-operator} rows and no party-local row. In the
pairing $\langle\lmoJ_t,S_t-J_t\rangle$ only the blocks of the gradient on
the joint sectors enter, $S_t$ and $J_t$ being lists over them. (H-SUPP) is
asked of $\cG^{J}_{\hat c,s}(J_s)$ on the whole sector $s$, and
$J_{(\mathsf t,s)}\succ0$ for every $\mathsf t$ suffices for it, since then
$J_s\succ0$. The weaker hypothesis per joint pair, which does not ask
anything of key-inert joint pairs inside a key-active sector, is
Corollary~\ref{cor:fo-dropped} at $\mathcal{D}=\varnothing$.

\emph{(2) The linear minimisation splits.} In trusted-source mode
$\Feas^{J}=\prod_{\mathsf t}\Feas^{J}_{\mathsf t}$
(Theorem~\ref{thm:choi}, general format), so
$\min_{J'\in\Feas^{J}}\langle\lmoJ_t,J'\rangle
=\sum_{\mathsf t}\min_{J'_{\mathsf t}\in\Feas^{J}_{\mathsf t}}
\langle\lmoJ_{t,\mathsf t},J'_{\mathsf t}\rangle$: one independent conic
program per Alice sector, each with its own coupling constraint of size
$d_{A_{0,\mathsf t}}$. A lower bound on each sector's minimum, summed over
the sectors, is therefore a lower bound on the minimum over the product,
which is how Lemma~\ref{lem:feasible-side-dual}\,(ii) forms the feasible
side of a bracket, one bound per Alice sector; and a pairing at most each
sector's minimum, summed, is at most the minimum over the product, which is
the second inequality of \eqref{eq:lmo}.
In untrusted-source mode the single normalisation and every scalar entry
couple all joint sectors and nothing splits. What does not split in either
mode is the iteration: the gradient couples all blocks through the scalars
$q(\hat c)/P_{\hat c}$ and $q(\hat c)/K_{\hat c}$.

\emph{(3) The relaxed feasible set per Alice sector.} In the general format
the \tref{relaxed feasible set} of Definition~\ref{def:relaxed-feasible-set}
and the recorded fields of Remarks~\ref{rem:fo-interface},
\ref{rem:fo-iterate-checks}, \ref{rem:fo-returned-point},
\ref{rem:fo-one-ruler}, \ref{rem:fo-reachability},
\ref{rem:fo-check-exposures}, \ref{rem:fo-oracle-repair},
\ref{rem:fo-iterate-repair} and~\ref{rem:fo-negative-values} are read with the
coupling constraint per Alice sector, and no statement of
Proposition~\ref{prop:relaxed-set-properties}, of
Remark~\ref{rem:feas-tol} or of those remarks changes;
what a run decides and records per sector is
Requirement~\ref{req:per-sector-coupling}.
\end{remark}
